% Annexe de la partie IS (refonte du 16/11/2026 ; coupée le 17/11/2026 selon synthese/CONSIGNES_COUPE.md : figures
% is_n02 et is_n06 reçues du corps ; tableaux d'hypothèses fusionnés). Chiffres : ../is/sorties/calculs_note_is.txt ;
% registre des hypothèses : ../revue/hypotheses_is.md ; tableaux générés par ../is/calculs_note_is.py (parties/tab_is_*.tex).

\subsection{Calibrage}\label{ann:is_calibrage}

Chiffres : \texttt{is/calculs\_note\_is.py}, sur la mécanique de \texttt{is/simulation\_is.py} corrigée le 16 novembre
2026, recalée sur 2025 ; euros courants, l'année 1 étant 2027. Faute de données de 2025, la comparaison de l'OCDE
(2023), les taux implicites de l'Insee (2022) et la répartition des bénéfices par taille (2024) restent antérieurs.

\begin{table}[!htbp]\centering\scriptsize
\caption{Données et paramètres de la partie IS}\label{tab:is_calibrage}
\begin{tabularx}{\linewidth}{@{}>{\raggedright\arraybackslash}p{4.6cm}>{\raggedright\arraybackslash}p{3.3cm}>{\raggedright\arraybackslash}X@{}}
\toprule
Grandeur & Valeur & Source ou statut \\
\midrule
Recette de 2025 (2,8 millions de redevables) & IS net 59,9~Md€ au budget ; en prélèvements obligatoires, 59,3 (69,5
avant 10,1 de crédits d'impôt), contribution exceptionnelle 7,5, contribution sociale 1,15 : 68,0~Md€ (2,27~\% du PIB) ;
minimum mondial : 0,5~Md€ en 2026 & PLF 2027, rapport AN n°~2847 ; Insee 3.217 \\
Investissement éligible des sociétés non financières (autres machines et équipements, TIC) & 84,1~Md€ (2025) ; 89,2~Md€
(2027) & Insee 8.212 ; +3~\% par an (3,2~\% de 2022 à 2025) \\
Stock correspondant ; usure & 541~Md€ (fin 2025) ; 14~\% par an & Insee 8.212 \\
Bénéfices imposés en 2025 & 30,6~Md€ à 15~\% ; 256,8~Md€ à 25~\% ; déficits 111,3~Md€ & DGFiP n°~51 (provisoire) \\
Croissance du bénéfice fiscal & +2~\% par an en 2026 et 2027 & hypothèse ; de 2022 à 2025, la base au taux normal a crû
de 2,4~\%, celle à 15~\% de 15~\% \\
Un point de taux normal en 2027 & 2,99~Md€ (2,67 si le taux réduit est gardé) & bénéfices de 2025 $\times$ 1,040 \\
Fin du taux réduit & 2,15 à 3,18~Md€ & haut : 10 points sur 30,6~Md€ ; bas : coût officiel de 2019 ramené à 10 points,
porté à 3~\% par an \\
CIR apprécié au niveau du groupe ; fin des zones hors outre-mer & 0,59 ; 0,54~Md€ (IR et IS) & CPO (2025) ; estimation
de 2021 vieillie ; part de l'IS non publiée \\
Boîte à brevets à 17,5~\% & 0,21 à 0,62~Md€ & moitié de la dépense fiscale (1,26~Md€ en 2025, estimée par le PLF 2026) \\
Ressources de l'assiette & 3,50 / 4,22 / 4,94~Md€ & bas, central, haut, en 2027, puis +3~\% par an \\
Durée d'amortissement des équipements & 7 ans, linéaire ou dégressif (coefficient 2,25) & BOFiP, BOI-BIC-AMT-20-20-30 \\
\bottomrule
\end{tabularx}
\end{table}

\figbloc{is_n02_qui_paie}{370 grandes entreprises paient 29~\% de l'IS, 2,6 millions de très petites
sociétés 21~\% (2025)}{fig:is_qui}

\textbf{Le plafond de 1~M€.} Parmi les 1,2 million d'entreprises qui ont investi en 2021, les 2~\% qui ont investi au
moins 500\,000~€ chacune réalisent 87~\% des 248,7~Md€ d'investissements corporels, et les 59~\% qui ont investi moins
de 10\,000~€ n'en font que 1~\% (Insee, 2023). Une interpolation de Pareto entre ces seuils, prolongée jusqu'à 1~M€,
donne une part couverte de 21~\% (18~\% pour un plafond de 0,5~M€). Le champ de l'Insee est plus large que celui de la
mesure, et l'on ignore la composition des achats par taille : nous retenons 20~\% (15 à 30~\%). Au
Royaume-Uni, 23~Md£ d'investissement ont été déduits en 2019-2020 sous le plafond équivalent (IFS, 2023).

\textbf{Expériences étrangères.} La déduction totale britannique exclut voitures et biens loués ; elle coûte de
l'ordre de 10 à 11~Md£ par an au début (IFS, 2023). La déduction américaine de 100~\% de 2017, réduite de
20 points par an dès 2023, a été rétablie pour les biens acquis après le 19 janvier 2025, occasion comprise sauf entre
entreprises liées (EY, 2026).

\subsection{Méthode}\label{ann:is_methode}

\textbf{Coût annuel.} Les achats de l'année $v$ sont $I_v = I_{2027}(1+g)^{v}$. Avec une fraction $b_v$ déduite
immédiatement et le reste amorti selon le plan actuel $d_k$, le coût pour l'État l'année $t$ est
$\tau\, s \sum_{v \le t} b_v I_v\,(\mathbf{1}_{t=v} - d_{t-v})$, où $\tau = 25~\%$ et $s$ est la part imposable. Le
plafond est une fraction $\theta$ des achats déduite à 100~\% : $b_v = \theta + (1-\theta)\,p_v$, où $p_v$ est le
palier de l'année. Le cas central est la moyenne des plans linéaire et dégressif.

\textbf{Coût actualisé et coût permanent.} Une année d'achats coûte $\tau\, s\, I\,(1 - z)$ en valeur actualisée ($z$ :
valeur actuelle des déductions pour 1~€ : 0,847 en linéaire, 0,897 en dégressif à 5~\% ; 0,896 et 0,930 à
3,25~\%). En régime permanent, le coût annuel vaut $\tau\, s\, I_t\,(1 - z(g))$, où $z(g)$ est actualisé au taux de
croissance de l'investissement : il serait nul sans croissance.

\textbf{Dette et baisse finale.} La dette vaut celle de l'année précédente, intérêts à 3,25~\% compris, moins le solde
de l'année. L'année qui suit l'extinction, le taux normal baisse de $x = (F - C)/(P - C/25)$ points, où $F$, $C$ et $P$
sont les ressources, le coût de la déduction et la valeur d'un point en régime permanent (2040) ; $C/25$ traduit la
baisse du coût de la déduction avec le taux.

\textbf{Taux effectifs.} Méthode de Devereux et Griffith (2003) : financement par fonds propres, sans impôt personnel,
rendement réel exigé de 5~\%, inflation de 2~\%, usure de 14~\% (équipements) ou 3,1~\% (bâtiments, amortis sur 25 ans) ;
taux moyen effectif pour un projet rapportant 20~\% (tableau~\ref{tab:is_tef}). L'OCDE (2025), dont les conventions
diffèrent (note de méthode non lue), donne pour la France un taux moyen effectif de 23,6~\% contre 22,2~\% en moyenne
(13\up{e} sur 38) et un taux marginal effectif de 15,0~\% contre 16,5~\% (9,8~\% contre 10,4~\% pour les équipements),
sans semble-t-il la contribution exceptionnelle. Seuls les écarts entre variantes se comparent.

\begin{table}[!htbp]\centering\scriptsize
\caption{Taux effectifs d'imposition d'un investissement, aujourd'hui et selon les réformes}\label{tab:is_tef}
\begin{tabularx}{\linewidth}{@{}>{\raggedright\arraybackslash}X>{\raggedleft\arraybackslash}p{1.4cm}>{\raggedleft\arraybackslash}p{2.2cm}>{\raggedleft\arraybackslash}p{2.6cm}@{}}
\toprule
Cas & Taux légal & Taux marginal effectif & Taux moyen effectif (projet à 20~\%) \\
\midrule
Équipement, aujourd'hui, amortissement linéaire & 25~\% & 20,7~\% & 23,6~\% \\
Équipement, aujourd'hui, amortissement dégressif & 25~\% & 15,1~\% & 22,1~\% \\
Équipement, baisse du taux de même coût actualisé que la déduction & 24,0~\% & 19,8~\% & 22,7~\% \\
Équipement, taux à 22~\% & 22~\% & 18,1~\% & 20,8~\% \\
Équipement, taux à 20~\% & 20~\% & 16,3~\% & 18,9~\% \\
Équipement, déduction immédiate de 50~\% & 25~\% & 11,5~\% & 21,2~\% \\
Équipement, déduction immédiate de 100~\% & 25~\% & 0 & 18,8~\% \\
Équipement, petite société sous 42\,500~€, aujourd'hui & 15~\% & 12,1~\% & 14,2~\% \\
Bâtiment, aujourd'hui & 25~\% & 22,0~\% & 24,0~\% \\
Bâtiment, petite société sous 42\,500~€, aujourd'hui & 15~\% & 13,0~\% & 14,4~\% \\
Bâtiment, taux à 22~\% & 22~\% & 19,2~\% & 21,1~\% \\
\bottomrule
\end{tabularx}
\par\vspace{3pt}\parbox{\linewidth}{\raggedright\scriptsize Coin sur le coût d'usage des équipements, au taux nominal de
5~\% : 5,1~\% en linéaire, 3,5~\% en dégressif, 0 avec la déduction immédiate. Source : \texttt{is/calculs\_note\_is.py}.}
\end{table}

\textbf{Part d'une baisse du taux qui touche l'investissement en équipements.} L'impôt actualisé sur le rendement normal
d'une année d'achats d'équipements égale le coût actualisé de leur déduction immédiate, 1,73 à 2,55~Md€. Un point de taux
en retire un vingt-cinquième, 0,07 à 0,10~Md€, pour un coût de 2,67~Md€ par an : 2,6 à 3,8~\% du coût, un peu plus si
l'on comptait les autres actifs neufs. Par taille, la baisse irait pour 30~\% aux grandes entreprises, 25~\% aux
entreprises de taille intermédiaire, 28 à 31~\% aux PME et 14 à 17~\% aux très petites sociétés (bénéfices de 2024).

\textbf{Production et bien-être.} Perte sèche : \texttt{synthese/sorties/perte\_seche\_autres.txt} ; élasticité :
H18-H19. Avec 50~\% sans montée, les gains de production et de bien-être sont divisés par deux ;
avec le plafond seul, le gain de production n'est que de 0,02~\% environ. Le suramortissement français de 2015 n'a pas
d'évaluation causale connue ; la seule étude du suramortissement robotique de 2019-2020 trouve un effet de 12 à 14~\%
sur l'ensemble des PME industrielles, difficile à concilier avec ses 2\,400 utilisateurs (DGE, 2026).

\textbf{Cas types.} Bénéfice et achats constants ; valeur de la déduction immédiate égale à l'avance d'impôt actualisée à
5~\% par année d'achats, $\tau\, I\,(1 - z)$, soit 3,8~\% des achats en linéaire ; avance de trésorerie de la première
année, $\tau\, I\,(1 - d_0)$ ; contribution exceptionnelle du grand groupe calculée sur un impôt supposé stable
(tableau~\ref{tab:is_cas}).

\begin{table}[!htbp]\centering\scriptsize
\caption{Quatre sociétés types : ce que change la proposition (€ par an)}\label{tab:is_cas}
\begin{tabularx}{\linewidth}{@{}>{\raggedright\arraybackslash}X>{\raggedleft\arraybackslash}p{1.45cm}>{\raggedleft\arraybackslash}p{1.45cm}>{\raggedleft\arraybackslash}p{1.55cm}>{\raggedleft\arraybackslash}p{1.9cm}>{\raggedleft\arraybackslash}p{1.4cm}@{}}
\toprule
Société & Impôt actuel & Fin du taux réduit & Valeur de la déduction & Variation nette & Trésorerie 2027 \\
\midrule
N'investit pas, bénéfice 42\,500~€ & 6\,375 & +4\,250 & 0 & +4\,250 (+67~\%) & +4\,250 \\
PME de services, bénéfice 40\,000~€, achats 10\,000~€ & 6\,000 & +4\,000 & $-$382 & +3\,618 (+60~\%) & +1\,679 \\
PME industrielle, bénéfice 300\,000~€, achats 150\,000~€ & 70\,750 & +4\,250 & $-$5\,727 & $-$1\,477 ($-$2~\%) & $-$30\,571 \\
Grand groupe, bénéfice 500~M€, achats 300~M€ & 129,1~M & 0 & $-$11,5~M & $-$11,5~M ($-$9~\%) & $-$34,9~M \\
Grand groupe : fin de la contribution exceptionnelle & 180,6~M (2026) & & & $-$51,5~M ($-$29~\%) & $-$51,5~M \\
\bottomrule
\end{tabularx}
\par\vspace{3pt}\parbox{\linewidth}{\raggedright\scriptsize Signe positif : la société paie davantage. Régime permanent,
avant la baisse finale du taux (3 à 6,5~M€ par an pour le grand groupe). Les PME achètent sous le plafond et déduisent
100~\% dès 2027 ; le grand groupe déduit 50~\% de ses achats au-delà de 1~M€ en 2027. Une société qui amortit
aujourd'hui en dégressif doit acheter 164\,000~€ par an pour compenser la perte de 4\,250~€.}
\end{table}

\subsection{Variantes}\label{ann:is_variantes}

\textbf{Calendriers.} Le tableau~\ref{tab:is_options} compare sept calendriers, la figure~\ref{fig:is_calendriers} le
coût annuel de trois d'entre eux, la figure~\ref{fig:is_dette_cal} la dette de ceux qui en créent une. L'amortissement
neutre coûte à l'État, à son taux, autant que la déduction immédiate (1,2 à 1,7~Md€ par année d'achats), mais ne vaut
pour les entreprises que 1,1 à 1,6~Md€, contre 1,7 à 2,6. Une marche unique en 2033 ferait bondir le coût de 2 à 10~Md€ ; le
gain d'attente qu'elle crée vaut un mois de rendement d'une machine rapportant 15 à 20~\% par an, contre une semaine
avec des paliers (H28). Le tableau~\ref{tab:is_fin} donne la dette selon le financement.

\begin{table}[!htbp]\centering\scriptsize
\caption{Les calendriers de la déduction immédiate comparés (Md€ ; part imposable 75~\%)}\label{tab:is_options}
\begin{tabularx}{\linewidth}{@{}>{\raggedright\arraybackslash}X>{\raggedleft\arraybackslash}p{1.45cm}>{\raggedleft\arraybackslash}p{1.45cm}>{\raggedleft\arraybackslash}p{1.45cm}>{\raggedleft\arraybackslash}p{1.45cm}>{\raggedleft\arraybackslash}p{1.45cm}@{}}
\toprule
Calendrier & Coût 2027 & Plus forte hausse annuelle & Coût actualisé par année d'achats & Dette maximale (cas central) & Extinction (central ; bornes) \\
\midrule
50~\% de 2027 à 2032, 100~\% en 2033 (ancienne note) & 7,0--7,8 & 7,9 & 1,7--2,6 & 11,2 & 2041 (2030 à 2051) \\
50~\% en 2027, +10 points par an, 100~\% en 2032 & 7,0--7,8 & 0,1 & 1,7--2,6 & 21,4 & 2042 (2036 à 2053) \\
100~\% sous 1 M€ dès 2027 ; au-delà, paliers 50 à 100~\% (2032) & 8,4--9,3 & 0,1 & 1,7--2,6 & 22,9 & 2042 (2036 à 2053) \\
100~\% sous 1 M€ par groupe, sans plus (AIA britannique seule) & 2,8--3,1 & 0,0 & 0,3--0,5 & aucune & --- \\
100~\% pour tous dès 2027 & 14,0--15,5 & 0,1 & 1,7--2,6 & 30,5 & 2043 (2036 à 2054) \\
50~\% pour tous, sans montée & 7,0--7,8 & 0,0 & 0,9--1,3 & 5,7 & 2033 (2030 à 2036) \\
Amortissement neutre (annuités revalorisées à 3,25~\%) & 0 & 0,4 & 1,1--1,6 & aucune & --- \\
\bottomrule
\end{tabularx}
\par\vspace{3pt}\parbox{\linewidth}{\raggedright\scriptsize Coût 2027 : régime actuel dégressif -- linéaire. Plus forte hausse annuelle : saut du coût d'une année sur l'autre (cas central) ; la falaise de 2033 en est le cas extrême. Coût actualisé à 5~\% (taux des entreprises), en régime permanent. Bornes de l'extinction : cas favorable (dégressif, ressources hautes) et défavorable (linéaire, ressources basses). Plafond : 1~M€ par groupe, couvrant 20~\% de l'investissement éligible (hypothèse). Amortissement neutre : plans actuels, annuités revalorisées au taux d'emprunt de l'État. Source : \texttt{is/calculs\_note\_is.py}.}
\end{table}


\figbloc{is_n06_calendriers}{Le calendrier proposé évite la marche de 2033 : son coût décroît dès 2027 au lieu de
bondir de 8~Md€}{fig:is_calendriers}

\figbloc{is_n09_dette_calendriers}{La dette de transition selon le calendrier : 29~Md€ au plus pour 100~\% dès 2027,
22~Md€ pour la proposition, 9~Md€ pour l'ancien calendrier}{fig:is_dette_cal}

\begin{table}[!htbp]\centering\scriptsize
\caption{La dette de transition selon le financement (proposition, cas central pour le coût)}\label{tab:is_fin}
\begin{tabularx}{\linewidth}{@{}>{\raggedright\arraybackslash}X rrr@{}}
\toprule
Financement & Ressources 2027 (Md€) & Dette maximale (Md€) & Extinction \\
\midrule
Assiette (proposition) & 3,5 à 4,9 & 16 à 28 & 2038 à 2046 \\
Assiette, taux réduit supprimé en deux ans (20~\% en 2027) & 2,9 en 2027 (central) & 23 (central) & 2041 (central) \\
Assiette, taux réduit porté à 20~\% & 2,4 à 3,4 & 29 à 40 & 2047 à 2067 \\
Assiette, taux réduit retiré aux seules sociétés immobilières et financières (hypothèse) & 1,9 à 2,7 & 36 (haut) ; 75 en
2066 (bas) & 2057 (haut) ; jamais (bas) \\
Assiette, taux réduit maintenu & 1,3 à 1,8 & 88 à 143 en 2066 & jamais \\
Taux normal à 28~\%, taux réduit gardé, sans autre mesure & 4,6 à 8,0 & 1 à 18 & 2030 à 2039 \\
Assiette et une dernière année de contribution exceptionnelle (5,0 en 2027, 0,5 en 2028) & 3,5 à 4,9 + 5,5 & 10 à 21
(central 15) & 2036 à 2043 (central 2039) \\
\bottomrule
\end{tabularx}
\par\vspace{3pt}\parbox{\linewidth}{\raggedright\scriptsize Régime actuel moitié linéaire, moitié dégressif ; part
imposable 75~\%. « Jamais » : dette non éteinte en 2066. Sociétés immobilières et financières : base à 15~\% supposée
répartie comme le nombre de sociétés (31~\% en 2025) ou l'IS brut (24~\%), invérifiable (H25). Source :
\texttt{is/calculs\_note\_is.py}.}
\end{table}

\subsection{La transition, année par année}\label{ann:is_transition}

\begin{table}[!htbp]\centering\scriptsize
\setlength{\tabcolsep}{3.5pt}
\caption{Transition année par année, avec les bornes (Md€ courants)}\label{tab:is_transition_complet}
\begin{tabular}{@{}l>{\raggedleft\arraybackslash}p{1.3cm}>{\raggedleft\arraybackslash}p{1.3cm}>{\raggedleft\arraybackslash}p{1.3cm}>{\raggedleft\arraybackslash}p{1.3cm}>{\raggedleft\arraybackslash}p{1.3cm}>{\raggedleft\arraybackslash}p{1.3cm}>{\raggedleft\arraybackslash}p{1.3cm}>{\raggedleft\arraybackslash}p{1.3cm}@{}}
\toprule
& \multicolumn{4}{c}{Cas central} & \multicolumn{2}{c}{Cas favorable} & \multicolumn{2}{c}{Cas défavorable}\\
\cmidrule(lr){2-5}\cmidrule(lr){6-7}\cmidrule(lr){8-9}
Année & Coût & Ressources & Solde & Dette & Coût & Dette & Coût & Dette \\
\midrule
2027 & 8,9 & 4,1 & $-$4,8 & 4,8 & 8,4 & 3,8 & 9,3 & 5,8 \\
2028 & 8,3 & 4,2 & $-$4,1 & 9,1 & 7,1 & 6,3 & 9,4 & 11,8 \\
2029 & 7,9 & 4,3 & $-$3,6 & 12,9 & 6,3 & 7,9 & 9,4 & 17,9 \\
2030 & 7,5 & 4,4 & $-$3,1 & 16,4 & 5,9 & 8,9 & 9,2 & 23,9 \\
2031 & 7,2 & 4,6 & $-$2,6 & 19,6 & 5,6 & 9,6 & 8,8 & 29,6 \\
2032 & 6,8 & 4,7 & $-$2,0 & 22,3 & 5,2 & 9,8 & 8,3 & 34,8 \\
2033 & 4,7 & 4,8 & 0,1 & 22,9 & 3,4 & 7,9 & 6,1 & 37,8 \\
2034 & 3,6 & 5,0 & 1,4 & 22,2 & 2,6 & 5,1 & 4,5 & 39,2 \\
2035 & 2,9 & 5,1 & 2,2 & 20,7 & 2,1 & 1,6 & 3,7 & 39,8 \\
2036 & 2,4 & 5,3 & 2,9 & 18,5 & 1,8 & $-$2,6 & 3,0 & 39,5 \\
2037 & 2,1 & 5,5 & 3,4 & 15,7 & 1,6 & $-$7,3 & 2,6 & 38,7 \\
2038 & 1,9 & 5,6 & 3,7 & 12,5 & 1,5 & $-$12,5 & 2,3 & 37,5 \\
2039 & 1,9 & 5,8 & 3,9 & 9,1 & 1,5 & $-$17,9 & 2,3 & 36,0 \\
2040 & 2,0 & 6,0 & 4,0 & 5,4 & 1,6 & $-$23,7 & 2,4 & 34,4 \\
2041 & 2,0 & 6,1 & 4,1 & 1,4 & 1,6 & $-$29,8 & 2,4 & 32,7 \\
2042 & 2,1 & 6,3 & 4,2 & $-$2,7 & 1,7 & $-$36,3 & 2,5 & 30,8 \\
2043 & 2,2 & 6,5 & 4,3 & $-$7,2 & 1,7 & $-$43,1 & 2,6 & 28,8 \\
2044 & 2,2 & 6,7 & 4,5 & $-$11,9 & 1,8 & $-$50,4 & 2,7 & 26,6 \\
2045 & 2,3 & 6,9 & 4,6 & $-$16,9 & 1,8 & $-$58,1 & 2,7 & 24,3 \\
2046 & 2,4 & 7,1 & 4,8 & $-$22,2 & 1,9 & $-$66,2 & 2,8 & 21,8 \\
2047 & 2,4 & 7,3 & 4,9 & $-$27,8 & 2,0 & $-$74,7 & 2,9 & 19,1 \\
2048 & 2,5 & 7,6 & 5,0 & $-$33,8 & 2,0 & $-$83,7 & 3,0 & 16,2 \\
2049 & 2,6 & 7,8 & 5,2 & $-$40,0 & 2,1 & $-$93,2 & 3,1 & 13,1 \\
2050 & 2,7 & 8,0 & 5,3 & $-$46,7 & 2,1 & $-$103,2 & 3,2 & 9,8 \\
2051 & 2,7 & 8,3 & 5,5 & $-$53,7 & 2,2 & $-$113,8 & 3,3 & 6,3 \\
2052 & 2,8 & 8,5 & 5,7 & $-$61,1 & 2,3 & $-$124,9 & 3,4 & 2,6 \\
2053 & 2,9 & 8,8 & 5,8 & $-$69,0 & 2,3 & $-$136,6 & 3,5 & $-$1,4 \\
2054 & 3,0 & 9,0 & 6,0 & $-$77,2 & 2,4 & $-$148,9 & 3,6 & $-$5,6 \\
2055 & 3,1 & 9,3 & 6,2 & $-$86,0 & 2,5 & $-$161,8 & 3,7 & $-$10,1 \\
2056 & 3,2 & 9,6 & 6,4 & $-$95,1 & 2,6 & $-$175,4 & 3,8 & $-$14,8 \\
\bottomrule
\end{tabular}
\par\vspace{3pt}\parbox{\linewidth}{\raggedright\scriptsize Cas favorable : régime actuel dégressif, ressources hautes (4,62~Md€ en 2027). Cas défavorable : régime actuel linéaire, ressources basses (3,50~Md€). Cas central : moitié-moitié, ressources au milieu (4,06~Md€). Une dette négative est un excédent cumulé : dans la proposition, il n'apparaît pas, car l'excédent est rendu par la baisse du taux normal dès l'année qui suit l'extinction (le tableau montre ce qui se passerait sans cette baisse). Part imposable : 75~\%. Source : \texttt{is/calculs\_note\_is.py}.}
\end{table}


\subsection{Hypothèses fortes et chiffres non vérifiés}\label{ann:is_hypotheses}

Registre : \texttt{revue/hypotheses\_is.md} (H1 à H22 ; H23 à H28 ajoutées ici). Effets sur la proposition, cas
central, sauf mention.

\begin{table}[!htbp]\centering\scriptsize
\caption{Hypothèses fortes de la partie IS}\label{tab:is_hyp1}
\begin{tabularx}{\linewidth}{@{}>{\raggedright\arraybackslash}p{4.4cm}>{\raggedright\arraybackslash}p{2.4cm}>{\raggedright\arraybackslash}X@{}}
\toprule
Hypothèse & Où elle joue & Ce qu'elle change si elle est fausse \\
\midrule
H1. 75~\% de l'investissement éligible fait par des sociétés imposables (60 à 90~\%), sans statistique publiée & coûts,
dette, PIB & L'hypothèse qui pèse le plus (§~\ref{sec:is_transition}). \\
H2. Les sociétés imposables utilisent toute la déduction l'année de l'achat, les autres ne coûtent rien & profil du
coût & Si les autres imputaient un jour tous leurs déficits : coût actualisé majoré d'un tiers au plus (0,6 à
0,85~Md€ par année d'achats). Revalorisation des déficits non chiffrable. \\
H3. Champ : machines hors transport et informatique des sociétés non financières & coûts & Avec le transport,
déduction totale à 23,3~Md€ la première année au lieu de 15,5 ; entreprises à l'IR : au plus 14,5~Md€ d'assiette. \\
H4. Amortissement type sur 7 ans & coût, coin & À 5 ou 10 ans, coût actualisé de 1,9 ou 3,5~Md€ au lieu de 2,6 ; à 10
ans, coin de 7,0~\% au lieu de 5,1. \\
H5. Plan actuel moitié linéaire, moitié dégressif (mélange réel non publié) & fourchettes & Dette de 13 à 32~Md€ avec
les ressources centrales. \\
H6. Investissement +3~\% par an & coût permanent, dette & À 2 ou 4~\% : dette de 20,2 ou 22,9~Md€. \\
H7. Bénéfice fiscal +2~\% par an en 2026 et 2027 & point de taux & De 0 à 2,5~\% : point de 2,87 à 3,02~Md€. \\
H8-H9. Actualisation à 5~\% pour les entreprises, emprunt de l'État à 3,25~\% & valeurs, seuils, dette & À 8~\%, la
déduction vaut plus pour les entreprises ; de 2,5 à 5~\%, dette de 20,9 à 23,1~Md€. \\
H10-H11. Aucun investissement suscité ; coûts comptés l'année de l'exercice & premières années, calendrier & Surcoût
temporaire, récupéré ensuite ; en caisse, un an plus tard. \\
H12. Bénéfices déclarés : 0 à 1,5~\% de moins par point de taux, sur tous les bénéfices & IS à 28~\%, baisse du taux &
Mesuré sur les multinationales : borne basse du 28~\% pessimiste. \\
H13. Coût réel du transfert de bénéfices : 0 à 100~\% de l'impôt perdu & perte sèche du 28~\% & Avec un pays d'accueil
à 12,5-15~\%, borne haute de 1,5 au lieu de 3,2~Md€. \\
H14. Fin du taux réduit : borne basse au coût officiel de 2019, non relevée au recalage (68~\% du rendement statique,
contre 75~\% auparavant) & ressources & Réactions (passage à l'IR, rémunération en
salaire) non chiffrées. \\
H15. Boîte à brevets : 0,62~Md€, 0,2 si la base fuit (borne arbitraire) & ressources & Quelques Md€ de dette. \\
H16-H17. CIR au niveau du groupe et zones : chiffrages du CPO & ressources & CIR estimé en 2021 (liasses de 2016) ;
part de l'IS des zones non publiée. Au plus 0,5~Md€ d'écart. \\
H18-H19. $r+\delta = 19~\%$ ; élasticité du stock d'équipements de 0,5 à 1,6 (0,5 sans source) & PIB, perte sèche &
Effet proportionnel ; dépend de la taille si la déduction agit par la trésorerie. \\
H20. Gain de PIB hors effet du financement & PIB & Voir §~\ref{sec:is_prop} (petites sociétés). \\
H21-H22. Taux de 25~\% pour tous dans le coût ; référence : droit en vigueur & coûts, recette & Négligeable ; fin de la
contribution exceptionnelle affichée (§~\ref{sec:is_arbitrages}). \\
H23. Plafond couvrant 20~\% de l'investissement éligible (15 à 30~\%) & premières années & Coût 2027 de 8,5 à
9,6~Md€, dette de 21,2 à 22,3~Md€. \\
H24. Ressources au milieu de la fourchette (4,22~Md€) & transition & Bornes de 3,5 et 4,9~Md€ : dette de 8 à 40~Md€
avec les bornes de H5. \\
H25. Taux réservé aux sociétés opérationnelles : base à 15~\% répartie comme le nombre de sociétés ou l'IS brut &
variante de financement & Rendement de 0,5 à 1,0~Md€ : ordre de grandeur indicatif, invérifiable. \\
H26. Taux effectifs : conventions de Devereux et Griffith, fonds propres & figures~\ref{fig:is_cout_egal} et
\ref{fig:is_avant_apres} & Avec emprunt, l'équipement est subventionné ; nous ne durcissons pas le plafonnement des
charges financières (3~M€ ou 30~\% du résultat avant intérêts et amortissements), qui épargne la plupart des PME. \\
H27. Cas types : déduction valorisée à 5~\% ; grand groupe sans CIR ni brevets & cas types & À 8-10~\%, la déduction
vaut 50 à 80~\% de plus pour les PME. \\
H28. Gain d'attente comparé à un rendement brut de 15 à 20~\% par an & calendrier & Ordre de grandeur, non mesuré. \\
\bottomrule
\end{tabularx}
\end{table}

\begin{table}[!htbp]\centering\scriptsize
\caption{Chiffres non vérifiés sur une source primaire (partie IS)}\label{tab:is_nonverif}
\begin{tabularx}{\linewidth}{@{}>{\raggedright\arraybackslash}p{5.6cm}>{\raggedright\arraybackslash}X@{}}
\toprule
Chiffre ou fait & Statut \\
\midrule
Environ 450 redevables de la contribution exceptionnelle & Source : le cours ; ni la DGFiP ni le PLF 2027 ne le donnent \\
Contribution exceptionnelle de 2026 : environ 7,0~Md€ & Déduite du PLF 2027 (5,0~Md€ en 2027, « en diminution de
2,0~Md€ ») \\
Suppression de l'article 15 du PLF 2027 en première lecture & Source secondaire (Légifiscal) \\
Bénéfices de 2025 ; croissance en 2026 et 2027 & DGFiP n°~51 provisoire ; tome I des Voies et moyens non publié \\
Contribution sociale de 2025 : 1,15~Md€ & Insee 3.217 ; le budget de l'État inscrit 1,5~Md€ en 2026, concept non rapproché \\
Contribution exceptionnelle dans les taux effectifs de l'OCDE & Déduite des données \\
Maffini, Xing et Devereux (2019) : +11~\% & Version de travail ; le résumé publié donne +2,1 à 2,5 points de taux
d'investissement \\
Bernini et Treibich (2016) ; Devereux, Liu et Loretz (2014) & Résumés seulement lus \\
Règles américaines (biens d'occasion, entreprises liées, 19 janvier 2025) & Source secondaire (EY, 2026) \\
CIR calculé société par société & D'après le CPO (2025) ; CGI non relu \\
Répartition de la base à 15~\% par type de société & Aucune statistique \\
\bottomrule
\end{tabularx}
\end{table}
